\documentclass[10pt]{article}
\usepackage[margin=20mm]{geometry}
\usepackage{amsmath,amssymb,mathtools,fontspec,unicode-math,xurl,hyperref}
\IfFontExistsTF{DejaVu Serif}{\setmainfont{DejaVu Serif}}{\setmainfont{Latin Modern Roman}}
\setmathfont{latinmodern-math.otf}
\setlength{\parindent}{0pt}\setlength{\parskip}{6pt}\sloppy
\title{Local distributions along exponential last-passage geodesics: a technical review}
\author{AI-assisted technical review}
\date{}
\begin{document}
\maketitle
Abstract: This review checks the explicit local distributions, their Palm normalization, and the final concentration-to-almost-sure-convergence implication in Martin, Sly and Zhang. One local error is established: the Bernoulli parameter in an alternative general-direction sampling representation is reversed. An exact correction is proved. The separate derivations of Propositions 1.5 and 1.6 agree with independent calculations. The complete geometric convergence argument is not certified.\par
\section*{Identity, version and scope}
The reviewed article is James B. Martin, Allan Sly and Lingfu Zhang, Convergence of the environment seen from geodesics in exponential last-passage percolation, J. Eur. Math. Soc. 27 (2025), 877–970, DOI 10.4171/JEMS/1594 (https://doi.org/10.4171/JEMS/1594) [main]. The exact publisher PDF is the primary source; PDF page \(j\) corresponds to printed page \(876+j\). The supplied author manuscript arXiv:2106.05242v2, dated 16 December 2024, is a distinct source [preprint]; its page 29 was compared only for the alternative representation.\par
Detailed checks cover the geometric convention in (2.4)–(2.5), Lemmas 4.1–4.2, the calculations underlying Propositions 1.5–1.6, and the implication from Propositions 10.1–10.2 to Theorems 1.1–1.2. The introduction’s dependency map and the block argument in Section 10 were examined at the level described below. The coupling, ergodicity and geometric estimates in Sections 2–3 and 5–9 were not all re-proved. No whole-paper correctness claim follows.\par
\section*{Definitions and checked dependencies}
Fix \(p=\rho\in(0,1)\) and write \(q=1-p\), \(a=pq\) and \(D=p^2+q^2\). We use the paper’s positive-integer geometric convention, \[\Pr(X=k)=p q^{k-1},\qquad k=1,2,\ldots.\] The two Palm components \(\Psi^{(1)}\) and \(\Psi^{(2)}\) have mixture weights \(q^2/D\) and \(p^2/D\) (printed p. 908). Under either component, the next right and left waiting times are independent. Lemma 4.2 (p. 910) gives their survival functions. This conditional independence is essential: mixing the marginals first and then multiplying them would be unjustified.\par
Equation reference: eq:survival. Multiplying the two survival functions separately within each component gives \[S_1(h)=(1+ph)(1+ah)e^{-h},\qquad
 S_2(h)=(1+qh)(1+ah)e^{-h},\quad h\geq0.\] Consequently \[
 \frac{q^2S_1(h)+p^2S_2(h)}D
 =\left(1+\frac{a h}D\right)(1+ah)e^{-h},\] which reproduces Proposition 1.5. Differentiating yields the Gamma\((1,1)\), Gamma\((2,1)\) and Gamma\((3,1)\) mixture weights \[\frac{p^4+q^4}{D},\qquad 2pq,\qquad \frac{2p^2q^2}{D}.\] They are nonnegative and sum to one. Their mean equals \(1/D\), consistent with the Palm jump intensity \(D\) and the length-bias normalization in Lemma 4.1 (pp. 908–909).\par
For the corner calculation on p. 911, the first conditional crossing probability is \[\int_0^\infty (1+ah)e^{-qh}p^2h e^{-ph}\,dh
 =p^2(1+2a).\] The other is \(q^2(1+2a)\). Mixing gives \(2p^2q^2(1+2a)/D\), as in Proposition 1.6; the symmetric value is \(3/8\). These checks do not use the erroneous alternative representation below.\par
\section*{Substantive finding}
Problem 1 (Reversed Bernoulli parameter in the alternative representation).  On printed p. 911, under An alternative representation of the weights on geodesics, the article asserts, in paraphrase, that a positive-integer Geometric\((p)\) variable satisfies \(X-1\overset d=BX\) with an independent Bernoulli\((p)\) variable. It then uses that parameter in the general-\(p\) representation of \(L(0,1)\) and of its minimum with \(L(1,0)\). The parameter must be \(q=1-p\).\par
\subsection*{Decisive check and original hypotheses}
All assumptions can be met with \(p=1/3\) and an independent \(B\sim\mathrm{Bernoulli}(p)\). Then \[\Pr(X-1=0)=p=\tfrac13,\qquad
 \Pr(BX=0)=1-p=\tfrac23.\] Both distributions are defined on the stated domains, so the claimed equality in distribution fails. An alternative convention does not resolve the discrepancy: equations (2.4)–(2.5) and their accompanying Bernoulli construction on p. 889 explicitly fix the convention.\par
Equation reference: eq:bern. The effect extends to the asserted continuous distribution. For independent \(E_3,E_4\sim\mathrm{Exp}(1)\) and \(B\sim\mathrm{Bernoulli}(b)\), \[
 \Pr\{(E_3+B E_4)/p>h\}=(1+bph)e^{-ph}.\] The printed choice \(b=p\) produces \((1+p^2h)e^{-ph}\), whereas Lemma 4.2 requires \((1+pqh)e^{-ph}\). Setting \(p=1/3\) and \(h=1\) gives respectively \(10e^{-1/3}/9\) and \(11e^{-1/3}/9\).\par
Let \(U=(E_1+E_2)/q\), independently of \(V=(E_3+BE_4)/p\). Then \[\Pr\{\min(U,V)>h\}=(1+qh)(1+bph)e^{-h}.\] The printed and correct survival probabilities at \(p=1/3,h=1\) are therefore \(50/(27e)\) and \(55/(27e)\) respectively. Thus the error changes an asserted sampling distribution; it is not merely a typographical difference with no mathematical consequence.\par
\subsection*{Correction and verification}
Replace the three general-\(p\) occurrences of Bernoulli\((p)\) in that paragraph by Bernoulli\((1-p)\). To prove the corrected geometric identity, let \(X'\) be an independent Geometric\((p)\) variable and \(B\sim\mathrm{Bernoulli}(q)\). Then \[\Pr(BX'=0)=p,\qquad
 \Pr(BX'=k)=q\,p q^{k-1}=p q^k=\Pr(X-1=k)\quad(k\geq1).\] Equation [eq:bern] (\#eq:bern) with \(b=q\) restores Lemma 4.2 and the displayed minimum representation. All stated independence assumptions and exponential scales remain unchanged.\par
\subsection*{Counterevidence, versions and impact}
When \(p=1/2\), the two parameters coincide, so the symmetric representation and its \((1/4,1/2,1/4)\) Gamma mixture remain valid. Equation [eq:survival] (\#eq:survival) and the corner calculation independently verify the principal explicit formulas through Lemma 4.2. The alternative paragraph is not an input to those preceding proofs or to the main convergence argument. The supported assessment is Minor, limited to this additional general-direction sampling representation; neither whole-paper failure nor a false convergence theorem is established.\par
The same parameter occurs on p. 29 of the supplied arXiv v2. This is a checked correspondence for this passage only, not an assertion of full version equivalence. A bounded title/erratum search found no separate correction; this does not establish priority or exclude unindexed author correspondence.\par
\section*{Convergence checks and limits}
For each fixed bounded cylinder test function and tolerance, the exponential bounds in Propositions 10.1–10.2 are summable. Borel–Cantelli, followed by a countable sequence of tolerances, gives almost sure convergence of that integral. A countable convergence-determining family in each finite coordinate projection yields finite-dimensional weak convergence on a common event. The limiting finite-dimensional probability measures provide tightness of the projections; coordinatewise compact sets with summable tail errors give tightness in the countable product. This justifies the final weak-convergence implication on p. 964, conditional on the concentration propositions.\par
The supporting block argument needs \(m>2s\) for independence within each parity class of blocks, and endpoint errors of order \(s/m\) when a local path environment is replaced by the corresponding full-path environment. These can be made small in the stated order of parameter choices. For integer increments \(d_i\) with \(\sum_i d_i^2\leq Dk\), a sufficient counting bound is \[\#\{(d_i)\}\leq 2^k
 \binom{\lfloor k\sqrt D\rfloor+k}{k}
 \leq [2e(1+\sqrt D)]^k.\] Combined with a sufficiently small bad-block probability this retains exponential decay. The factor-two radius notation on p. 966 is not used here as a separate substantive finding: a path with \(2km+1\) vertices corresponds to radius \(km\) under the definition on p. 881. These local checks do not replace a complete verification of Proposition 8.1, Lemma 10.3, the Busemann estimates, or the integrable-probability inputs.\par
\section*{Conclusion}
Dimension | State\par
Supported issue | Minor: reversed parameter in an alternative general-direction representation.\par
Repair verification | Analytic distributional proof, valid for all \(p\in(0,1)\).\par
Surviving checked claims | Lemma 4.2-based formulas and the final conditional Borel–Cantelli implication.\par
Author adoption | Not requested; unknown.\par
Scientific decision | Pending; no human review supplied.\par
Remaining work | Full verification of the long geometric and coupling dependencies.\par
\section*{Supporting checks}
00000TOKEN\par
All displayed counterexample values are exact analytic calculations, not empirical experiments. Two native in-conversation helpers supplied bounded cross-checks of the waiting-time and concentration arguments; the primary executor evaluated the conclusions and performed the source-image and report-image checks. This is initial AI review, not independent acceptance. The controller’s UI observation and the executor’s runtime self-description are separately recorded; backend equivalence is not asserted.\par
No author code was executed. The acquisition record’s bounded supplement/code search found no dedicated release for this paper. That absence is a search limitation, not a defect.\par
9 J. B. Martin, A. Sly and L. Zhang, Convergence of the environment seen from geodesics in exponential last-passage percolation, J. Eur. Math. Soc. 27 (2025), 877–970. doi:10.4171/JEMS/1594 (https://doi.org/10.4171/JEMS/1594). Exact supplied publisher PDF. J. B. Martin, A. Sly and L. Zhang, same title, arXiv:2106.05242v2 (https://arxiv.org/abs/2106.05242v2), 16 December 2024. Alternative-representation passage on p. 29 only.\par
\end{document}
